% ESTILO DE MEMORIA ACADÉMICA
% Adaptación visual del TFG del autor: https://github.com/nowelito28/TFG
% La portada utiliza la identidad UCM y el contenido de esta práctica.
\definecolor{Granate}{HTML}{871B36}
\definecolor{Grafito}{HTML}{202020}
\definecolor{GrisTexto}{HTML}{555555}
\definecolor{GrisFondo}{HTML}{F8F8F8}
\definecolor{GrisLinea}{HTML}{B5B5B5}

\setstretch{1.35}
\setlength{\parindent}{1.5em}
\setlength{\parskip}{5pt plus 1pt minus 1pt}
\setlength{\emergencystretch}{2em}
\setcounter{tocdepth}{1}
\setcounter{secnumdepth}{2}
\setlist[itemize]{leftmargin=*,itemsep=2pt,topsep=3pt}
\setlist[enumerate]{leftmargin=*,itemsep=2pt,topsep=3pt}

% Capítulos como en el TFG: regla fina, etiqueta y título a la izquierda,
% y regla gruesa de cierre. Sin dependencias de fncychap.
\newcommand{\ReglaCapitulo}{\titlerule[0.5mm]}
\titleformat{\chapter}[display]
  {\normalfont\bfseries\Huge}
  {\titlerule\filright\Large\chaptertitlename\ \thechapter}
  {0mm}{\filright}
  [\vspace{0.5mm}\ReglaCapitulo]
\titleformat{name=\chapter,numberless}[display]
  {\normalfont\bfseries\Huge}
  {\titlerule}
  {0mm}{\filright}
  [\vspace{0.5mm}\ReglaCapitulo]
\titlespacing*{\chapter}{0pt}{35pt}{28pt}
\titleformat{\section}{\normalfont\Large\bfseries}{\thesection}{0.7em}{}
\titleformat{\subsection}{\normalfont\large\bfseries}{\thesubsection}{0.7em}{}
\titlespacing*{\section}{0pt}{16pt}{7pt}
\titlespacing*{\subsection}{0pt}{12pt}{5pt}

% La apertura lleva número al pie; las demás páginas, cabecera de capítulo.
\pagestyle{fancy}
\fancyhf{}
\fancyhead[L]{\small\itshape\nouppercase{\leftmark}}
\fancyhead[R]{\small\thepage}
\renewcommand{\headrulewidth}{0pt}
\renewcommand{\footrulewidth}{0pt}
\renewcommand{\chaptermark}[1]{%
  \markboth{\ifnum\value{chapter}>0 \chaptername\ \thechapter.\ \fi #1}{}%
}
\fancypagestyle{plain}{%
  \fancyhf{}
  \fancyfoot[C]{\thepage}
  \renewcommand{\headrulewidth}{0pt}
}

\hypersetup{
  unicode=true,hidelinks,
  pdftitle={\TituloMemoria: \Modulo},
  pdfauthor={\Autor},
  pdfsubject={\Titulacion},
  pdfkeywords={criptografía, esteganografía, memoria, práctica},
  bookmarksnumbered=true
}
\urlstyle{same}
\captionsetup{font=small,labelfont=bf,labelsep=period,skip=7pt,justification=centering}
\addto\captionsspanish{\renewcommand{\contentsname}{Índice general}}

% Marco discreto para distinguir fielmente cada enunciado del desarrollo.
\tcbset{memoria/.style={
  enhanced,sharp corners,boxrule=0.4pt,
  colframe=GrisLinea,colback=white,
  left=10pt,right=10pt,top=7pt,bottom=7pt,
  before skip=9pt,after skip=9pt,
  fonttitle=\normalfont\small\bfseries,
  coltitle=black,colbacktitle=GrisFondo,
  before upper={\setstretch{1.15}\setlength{\parindent}{0pt}}
}}
\newtcolorbox{enunciado}[1]{memoria,title={Enunciado #1}}
\newtcolorbox{notaestado}{memoria,title={Estado del documento}}

% Los futuros listados tienen marco superior e inferior, como en el TFG.
\lstdefinestyle{terminal}{
  basicstyle=\ttfamily\small,
  backgroundcolor=\color{white},
  frame=tb,rulecolor=\color{GrisLinea},
  breaklines=true,columns=fullflexible,keepspaces=true,
  showstringspaces=false,tabsize=4,
  aboveskip=3mm,belowskip=3mm,
  literate={á}{{\'a}}1 {é}{{\'e}}1 {í}{{\'i}}1 {ó}{{\'o}}1 {ú}{{\'u}}1
           {ñ}{{\~n}}1 {Á}{{\'A}}1 {É}{{\'E}}1 {Í}{{\'I}}1 {Ó}{{\'O}}1 {Ú}{{\'U}}1
}
\lstset{style=terminal}
\renewcommand{\lstlistingname}{Código}
\lstdefinestyle{programacompleto}{
  style=terminal,basicstyle=\ttfamily\footnotesize,
  numbers=left,numberstyle=\tiny\color{GrisTexto},
  numbersep=8pt,xleftmargin=6mm
}
